% ===========================================================================
\chapter{Especificação dos Casos de Uso de Auditoria Clínica e Ferramentas da Qualidade}
\label{cap_auditoria_gtt}
% ===========================================================================

O núcleo operacional e pedagógico do SIGEA-GTT é a realização da auditoria retrospectiva simulada pelo discente. Neste subsistema, o estudante é imerso em uma interface de alta fidelidade clínica para analisar o prontuário fictício sob a pressão temporal de vinte minutos estipulada pelo protocolo internacional do IHI \cite{ihi2009}. Ao detectar rastros de eventos adversos, o discente aciona os gatilhos correspondentes e constrói o plano de melhoria contínua através das quatro ferramentas da qualidade.

A Figura \ref{fig_uc_modulo_auditoria} ilustra o diagrama de casos de uso deste subsistema em notação UML com renderização em TikZ:

\begin{figure}[!ht]
\centering
\caption{Diagrama de Casos de Uso do Subsistema de Auditoria e Qualidade}
\label{fig_uc_modulo_auditoria}
\resizebox{0.92\textwidth}{!}{%
\begin{tikzpicture}[
    actor/.style={
        draw=clinical-900,
        fill=clinical-100!40,
        thick,
        rounded corners=2pt,
        align=center,
        font=\bfseries\scriptsize,
        minimum width=2.4cm,
        minimum height=1.2cm
    },
    usecase/.style={
        ellipse,
        draw=clinical-700,
        fill=white,
        thick,
        align=center,
        font=\scriptsize,
        minimum width=3.4cm,
        minimum height=1.1cm,
        inner sep=2pt
    },
    usecase_tool/.style={
        ellipse,
        draw=teal!80!black,
        fill=teal!5,
        thick,
        align=center,
        font=\scriptsize,
        minimum width=3.2cm,
        minimum height=1.0cm,
        inner sep=2pt
    },
    link/.style={
        draw=clinical-900!80,
        thick
    },
    include_link/.style={
        draw=clinical-500,
        thick,
        dashed,
        -Latex
    }
]
    % Fronteira
    \draw[rounded corners=8pt, fill=clinical-100!15, draw=clinical-700, line width=1.2pt] 
        (0.5, -9.5) rectangle (14.2, 3.2);
    \node[anchor=north west, font=\bfseries\small\color{clinical-700}] at (0.8, 2.9) 
        {Subsistema: Auditoria GTT e Ferramentas da Qualidade};

    % Atores
    \node[actor] (student) at (-2.2, -1.0) {Discente Auditor\\\texttt{\scriptsize (STUDENT)}};
    \node[actor, draw=alert-amber!90!black, fill=alert-amber!10] (timer) at (16.8, 1.5) {Ator Temporal\\\texttt{\scriptsize (SYSTEM\_TIMER)}};

    % Casos de Uso de Auditoria
    \node[usecase] (uc11) at (3.5, 1.5) {UC11: Iniciar Sessão\\de Auditoria};
    \node[usecase] (uc12) at (8.5, 1.5) {UC12: Controlar Cronômetro\\de 20 Minutos};
    \node[usecase] (uc13) at (3.5, -1.2) {UC13: Rastrear e Vincular\\Gatilhos no Prontuário};

    % Casos de Uso das Ferramentas
    \node[usecase_tool] (uc14) at (3.5, -4.2) {UC14: Preencher\\Ishikawa (6M)};
    \node[usecase_tool] (uc15) at (8.5, -4.2) {UC15: Construir\\Matriz GUT};
    \node[usecase_tool] (uc16) at (3.5, -6.5) {UC16: Elaborar\\Plano 5W2H};
    \node[usecase_tool] (uc17) at (8.5, -6.5) {UC17: Estruturar\\Ciclo PDCA};

    % Submissão
    \node[usecase] (uc18) at (6.0, -8.5) {UC18: Submeter Auditoria\\Clínica Completa};

    % Conexões Ator Aluno
    \draw[link] (student) -- (uc11);
    \draw[link] (student) -- (uc13);
    \draw[link] (student) -- (uc14);
    \draw[link] (student) -- (uc15);
    \draw[link] (student) -- (uc16);
    \draw[link] (student) -- (uc17);
    \draw[link] (student) -- (uc18);

    % Conexão Ator Timer
    \draw[link, draw=alert-amber!90!black, thick] (timer) -- (uc12);

    % Relações Include
    \draw[include_link] (uc11) -- node[above, font=\tiny\bfseries] {<<include>>} (uc12);
    \draw[include_link] (uc18) -- node[left, font=\tiny\bfseries] {<<include>>} (uc13);
    \draw[include_link] (uc18) -- (uc14);
    \draw[include_link] (uc18) -- (uc15);
    \draw[include_link] (uc18) -- (uc16);
    \draw[include_link] (uc18) -- (uc17);

\end{tikzpicture}
}
\fonte{Elaborado pelo autor (2026).}
\end{figure}

% ===========================================================================
\section{Especificação Detalhada dos Casos de Uso}
\label{sec_especificacao_uc_auditoria}
% ===========================================================================

% ---------------------------------------------------------------------------
% UC11
% ---------------------------------------------------------------------------
\subsection{UC11: Iniciar Sessão de Auditoria Prática}
\label{uc11_especificacao}

\begin{table}[!ht]
\centering
\caption{Especificação Detalhada do Caso de Uso UC11}
\label{tab_uc11}
\footnotesize
\begin{tabularx}{\textwidth}{p{3.5cm} X}
\toprule
\textbf{Propriedade} & \textbf{Especificação Técnica e Comportamental} \\
\midrule
\textbf{Identificador e Nome} & \texttt{UC11: Iniciar Sessão de Auditoria Prática} \\
\textbf{Atores} & \texttt{STUDENT} (Ator Primário). \\
\textbf{Nível de Meta} & Nível Usuário (\textit{Sea Level}). \\
\textbf{Pré-condições} & Discente matriculado na turma; atividade aberta e dentro do prazo limite (\texttt{deadline}); ausência de submissão prévia da mesma atividade pelo aluno. \\
\textbf{Gatilho} & O estudante clica no botão "Iniciar Auditoria" no cartão da atividade. \\
\textbf{Fluxo Principal} & 
1. O sistema verifica a elegibilidade do aluno e a validade do prazo.\\
& 2. O sistema exibe modal de instruções e termos de responsabilidade do método GTT.\\
& 3. O estudante confirma o início da auditoria.\\
& 4. O sistema carrega o caso clínico estruturado da coluna \texttt{clinical\_case\_data}.\\
& 5. O sistema inicializa o painel de auditoria dividindo a tela entre o prontuário e as ferramentas de anotação.\\
& 6. O sistema invoca compulsoriamente \texttt{UC12: Controlar Cronômetro de Auditoria}, disparando a contagem regressiva a partir de 20 minutos (1.200 segundos). \\
\textbf{Fluxos de Exceção} & 
\textbf{[1e] Prazo Limite Expirado}: O sistema bloqueia o início e informa que a data de entrega foi encerrada.\\
& \textbf{[1e2] Atividade Já Submetida}: A restrição de unicidade impede uma segunda tentativa, exibindo a nota ou status de avaliação. \\
\textbf{Pós-condições} & Sessão de auditoria iniciada e cronômetro em execução. \\
\textbf{Regras de Negócio} & Submissão única por discente e auditoria sob cronômetro. \\
\textbf{Requisitos Associados} & \texttt{RF14}, \texttt{RF15}, \texttt{RNF03}. \\
\bottomrule
\end{tabularx}
\fonte{Elaborado pelo autor (2026).}
\end{table}

% ---------------------------------------------------------------------------
% UC12
% ---------------------------------------------------------------------------
\subsection{UC12: Controlar Cronômetro de Auditoria (20 Minutos)}
\label{uc12_especificacao}

\begin{table}[!ht]
\centering
\caption{Especificação Detalhada do Caso de Uso UC12}
\label{tab_uc12}
\footnotesize
\begin{tabularx}{\textwidth}{p{3.5cm} X}
\toprule
\textbf{Propriedade} & \textbf{Especificação Técnica e Comportamental} \\
\midrule
\textbf{Identificador e Nome} & \texttt{UC12: Controlar Cronômetro de Auditoria (20 Minutos)} \\
\textbf{Atores} & \texttt{SYSTEM\_TIMER} (Ator Primário), \texttt{STUDENT} (Ator Secundário). \\
\textbf{Nível de Meta} & Nível Subfunção / Apoio Temporal (\textit{Fish Level}). \\
\textbf{Pré-condições} & Sessão de auditoria iniciada com sucesso em \texttt{UC11}. \\
\textbf{Gatilho} & Disparo pelo caso de uso \texttt{UC11} na abertura da interface de auditoria. \\
\textbf{Fluxo Principal} & 
1. O cronômetro regressivo é exibido no topo da tela em destaque visual.\\
& 2. O tempo decresce segundo a segundo de forma síncrona.\\
& 3. Ao atingir o minuto 15 (restam 5 minutos), o sistema emite um alerta sonoro suave e altera a cor do contador para amarelo âmbar (\texttt{\#D97706}).\\
& 4. Ao atingir o minuto 18 (restam 2 minutos), o sistema emite aviso sonoro e altera a cor para carmim de urgência (\texttt{\#DC2626}).\\
& 5. Ao atingir zero (vigésimo minuto), o sistema congela a interface de edição.\\
& 6. O sistema aciona automaticamente a submissão forçada dos dados coletados até aquele instante, invocando o fechamento de \texttt{UC18}. \\
\textbf{Fluxos Alternativos} & 
\textbf{[2a] Pausa Operacional}: O discente aciona o botão "Pausar Cronômetro" para tirar dúvidas didáticas com o docente. O cronômetro congela a contagem e oculta temporariamente o prontuário para evitar leitura em tempo pausado. O estudante retoma a contagem ao despausar. \\
\textbf{Pós-condições} & Tempo controlado e submissão garantida dentro da janela de 20 minutos. \\
\textbf{Regras de Negócio} & \textbf{RN09} (Limite Estrito de 20 Minutos conforme Protocolo IHI). \\
\textbf{Requisitos Associados} & \texttt{RF15}, \texttt{RNF03}. \\
\bottomrule
\end{tabularx}
\fonte{Elaborado pelo autor (2026).}
\end{table}

% ---------------------------------------------------------------------------
% UC13
% ---------------------------------------------------------------------------
\subsection{UC13: Rastrear e Vincular Gatilhos no Prontuário}
\label{uc13_especificacao}

\begin{table}[!ht]
\centering
\caption{Especificação Detalhada do Caso de Uso UC13}
\label{tab_uc13}
\footnotesize
\begin{tabularx}{\textwidth}{p{3.5cm} X}
\toprule
\textbf{Propriedade} & \textbf{Especificação Técnica e Comportamental} \\
\midrule
\textbf{Identificador e Nome} & \texttt{UC13: Rastrear e Vincular Gatilhos no Prontuário} \\
\textbf{Atores} & \texttt{STUDENT} (Ator Primário). \\
\textbf{Nível de Meta} & Nível Usuário (\textit{Sea Level}). \\
\textbf{Pré-condições} & Sessão de auditoria ativa com prontuário visível. \\
\textbf{Gatilho} & O estudante identifica uma pista clínica suspeita no prontuário. \\
\textbf{Fluxo Principal} & 
1. O estudante seleciona o trecho textual do prontuário que evidencia o evento.\\
& 2. O estudante clica em "Adicionar Gatilho".\\
& 3. O sistema abre o seletor com os 53 gatilhos canônicos agrupados por módulo.\\
& 4. O discente escolhe o gatilho aplicável (ex.: \texttt{"M03 - Anticoagulantes"}).\\
& 5. O sistema preenche automaticamente o campo de citação com o trecho selecionado.\\
& 6. O estudante informa a seção do prontuário (ex.: \textit{"Prescrição Médica"}).\\
& 7. O estudante digita a justificativa clínica formulando a hipótese de dano.\\
& 8. O estudante classifica a gravidade na escala NCC MERP (categorias A a I).\\
& 9. O sistema assinala automaticamente se houve dano real (\texttt{is\_harm = true} para E até I e \texttt{false} para A até D).\\
& 10. O gatilho rastreado é adicionado à lista lateral de achados da auditoria. \\
\textbf{Fluxos de Exceção} & 
\textbf{[7e] Justificativa Vazia}: O sistema exige fundamentação textual antes de salvar.\\
& \textbf{[8e] Ausência de Gravidade NCC MERP}: A seleção de gravidade é obrigatória. \\
\textbf{Pós-condições} & Gatilho registrado no array de achados temporários da sessão discente. \\
\textbf{Regras de Negócio} & \textbf{RN10} (Citação Textual e Justificativa Obrigatórias) e \textbf{RN11} (Escala NCC MERP). \\
\textbf{Requisitos Associados} & \texttt{RF16}, \texttt{RNF03}. \\
\bottomrule
\end{tabularx}
\fonte{Elaborado pelo autor (2026).}
\end{table}

% ---------------------------------------------------------------------------
% UC14
% ---------------------------------------------------------------------------
\subsection{UC14: Preencher Diagrama de Ishikawa (6M)}
\label{uc14_especificacao}

\begin{table}[!ht]
\centering
\caption{Especificação Detalhada do Caso de Uso UC14}
\label{tab_uc14}
\footnotesize
\begin{tabularx}{\textwidth}{p{3.5cm} X}
\toprule
\textbf{Propriedade} & \textbf{Especificação Técnica e Comportamental} \\
\midrule
\textbf{Identificador e Nome} & \texttt{UC14: Preencher Diagrama de Ishikawa (6M)} \\
\textbf{Atores} & \texttt{STUDENT} (Ator Primário). \\
\textbf{Nível de Meta} & Nível Usuário (\textit{Sea Level}). \\
\textbf{Pré-condições} & Pelo menos um gatilho clínico com dano identificado em \texttt{UC13}. \\
\textbf{Gatilho} & O estudante navega para a aba "Ishikawa" na área de análise causal. \\
\textbf{Fluxo Principal} & 
1. O sistema renderiza o espinha de peixe interativo com seis ramificações.\\
& 2. O estudante define o problema central investigado (cabeça do peixe).\\
& 3. O estudante insere causas primárias e secundárias nas dimensões aplicáveis: Método, Máquina, Material, Mão de Obra, Medida e Meio Ambiente.\\
& 4. O sistema atualiza o diagrama visual vetorial dinamicamente.\\
& 5. O sistema persiste a estrutura no nó \texttt{ishikawa} dos dados temporários. \\
\textbf{Fluxos de Exceção} & 
\textbf{[2e] Efeito Central Não Definido}: O sistema sinaliza a pendência no diagrama. \\
\textbf{Pós-condições} & Diagrama de causa-efeito 6M registrado para a submissão. \\
\textbf{Regras de Negócio} & Abordagem sistemática de análise de causas de incidentes hospitalares. \\
\textbf{Requisitos Associados} & \texttt{RF17}, \texttt{RNF04}. \\
\bottomrule
\end{tabularx}
\fonte{Elaborado pelo autor (2026).}
\end{table}

% ---------------------------------------------------------------------------
% UC15
% ---------------------------------------------------------------------------
\subsection{UC15: Construir Matriz de Priorização GUT}
\label{uc15_especificacao}

\begin{table}[!ht]
\centering
\caption{Especificação Detalhada do Caso de Uso UC15}
\label{tab_uc15}
\footnotesize
\begin{tabularx}{\textwidth}{p{3.5cm} X}
\toprule
\textbf{Propriedade} & \textbf{Especificação Técnica e Comportamental} \\
\midrule
\textbf{Identificador e Nome} & \texttt{UC15: Construir Matriz de Priorização GUT} \\
\textbf{Atores} & \texttt{STUDENT} (Ator Primário). \\
\textbf{Nível de Meta} & Nível Usuário (\textit{Sea Level}). \\
\textbf{Pré-condições} & Sessão de auditoria ativa com causas identificadas no Ishikawa. \\
\textbf{Gatilho} & O estudante acessa a aba "Matriz GUT". \\
\textbf{Fluxo Principal} & 
1. O sistema exibe a tabela de priorização com campos de pontuação.\\
& 2. O estudante adiciona uma linha descrevendo a fragilidade clínica encontrada.\\
& 3. O estudante seleciona a nota de Gravidade ($G$) de 1 a 5.\\
& 4. O estudante seleciona a nota de Urgência ($U$) de 1 a 5.\\
& 5. O estudante seleciona a nota de Tendência ($T$) de 1 a 5.\\
& 6. O sistema calcula automaticamente o score final: $GUT = G \times U \times T$.\\
& 7. O sistema reordena a tabela de forma decrescente pela pontuação final ($[1, 125]$). \\
\textbf{Fluxos de Exceção} & 
\textbf{[3e] Nota Fora do Intervalo}: O componente restringe a seleção estritamente ao conjunto discreto $\{1, 2, 3, 4, 5\}$. \\
\textbf{Pós-condições} & Matriz GUT calculada e ordenada por criticidade de intervenção. \\
\textbf{Regras de Negócio} & Cálculo exato $GUT = G \times U \times T$. \\
\textbf{Requisitos Associados} & \texttt{RF18}, \texttt{RNF03}. \\
\bottomrule
\end{tabularx}
\fonte{Elaborado pelo autor (2026).}
\end{table}

% ---------------------------------------------------------------------------
% UC16
% ---------------------------------------------------------------------------
\subsection{UC16: Elaborar Plano de Ação 5W2H}
\label{uc16_especificacao}

\begin{table}[!ht]
\centering
\caption{Especificação Detalhada do Caso de Uso UC16}
\label{tab_uc16}
\footnotesize
\begin{tabularx}{\textwidth}{p{3.5cm} X}
\toprule
\textbf{Propriedade} & \textbf{Especificação Técnica e Comportamental} \\
\midrule
\textbf{Identificador e Nome} & \texttt{UC16: Elaborar Plano de Ação 5W2H} \\
\textbf{Atores} & \texttt{STUDENT} (Ator Primário). \\
\textbf{Nível de Meta} & Nível Usuário (\textit{Sea Level}). \\
\textbf{Pré-condições} & Causas priorizadas na Matriz GUT em \texttt{UC15}. \\
\textbf{Gatilho} & O discente navega para a aba "Plano 5W2H". \\
\textbf{Fluxo Principal} & 
1. O sistema exibe o quadro operacional das sete dimensões diretivas.\\
& 2. O discente cadastra uma ação preventiva/corretiva preenchendo: O que (\textit{What}), Por que (\textit{Why}), Onde (\textit{Where}), Quando (\textit{When}), Quem (\textit{Who}), Como (\textit{How}) e Quanto custa (\textit{How Much}).\\
& 3. O sistema valida que nenhuma dimensão estratégica ficou em branco.\\
& 4. O sistema insere a ação na coleção estruturada \texttt{plan\_5w2h}. \\
\textbf{Fluxos de Exceção} & 
\textbf{[3e] Dimensão Não Preenchida}: O sistema sinaliza o campo faltante com borda vermelha e impede o avanço. \\
\textbf{Pós-condições} & Plano executivo 5W2H configurado e validado. \\
\textbf{Regras de Negócio} & Completude dos 7 campos operacionais. \\
\textbf{Requisitos Associados} & \texttt{RF19}, \texttt{RNF04}. \\
\bottomrule
\end{tabularx}
\fonte{Elaborado pelo autor (2026).}
\end{table}

% ---------------------------------------------------------------------------
% UC17
% ---------------------------------------------------------------------------
\subsection{UC17: Estruturar Ciclo PDCA}
\label{uc17_especificacao}

\begin{table}[!ht]
\centering
\caption{Especificação Detalhada do Caso de Uso UC17}
\label{tab_uc17}
\footnotesize
\begin{tabularx}{\textwidth}{p{3.5cm} X}
\toprule
\textbf{Propriedade} & \textbf{Especificação Técnica e Comportamental} \\
\midrule
\textbf{Identificador e Nome} & \texttt{UC17: Estruturar Ciclo PDCA} \\
\textbf{Atores} & \texttt{STUDENT} (Ator Primário). \\
\textbf{Nível de Meta} & Nível Usuário (\textit{Sea Level}). \\
\textbf{Pré-condições} & Plano de ação formulado na etapa 5W2H. \\
\textbf{Gatilho} & O estudante clica na aba "Ciclo PDCA". \\
\textbf{Fluxo Principal} & 
1. O sistema apresenta os quatro quadrantes do ciclo de Deming: Planejar (\textit{Plan}), Executar (\textit{Do}), Checar (\textit{Check}) e Agir (\textit{Act}).\\
& 2. O discente formula a meta de melhoria no quadrante \textit{Plan}.\\
& 3. O discente delineia as rotinas assistenciais no quadrante \textit{Do}.\\
& 4. O discente estabelece indicadores de checagem no quadrante \textit{Check}.\\
& 5. O discente define os padrões corretivos e de padronização em \textit{Act}.\\
& 6. O sistema consolida o objeto \texttt{cycle\_pdca} para envio. \\
\textbf{Pós-condições} & Estrutura de melhoria contínua PDCA finalizada. \\
\textbf{Regras de Negócio} & Fechamento do ciclo causal de segurança do paciente. \\
\textbf{Requisitos Associados} & \texttt{RF20}, \texttt{RNF04}. \\
\bottomrule
\end{tabularx}
\fonte{Elaborado pelo autor (2026).}
\end{table}

% ---------------------------------------------------------------------------
% UC18
% ---------------------------------------------------------------------------
\subsection{UC18: Submeter Auditoria Clínica Completa}
\label{uc18_especificacao}

\begin{table}[!ht]
\centering
\caption{Especificação Detalhada do Caso de Uso UC18}
\label{tab_uc18}
\footnotesize
\begin{tabularx}{\textwidth}{p{3.5cm} X}
\toprule
\textbf{Propriedade} & \textbf{Especificação Técnica e Comportamental} \\
\midrule
\textbf{Identificador e Nome} & \texttt{UC18: Submeter Auditoria Clínica Completa} \\
\textbf{Atores} & \texttt{STUDENT} (Ator Primário), \texttt{SYSTEM\_TIMER} (Ator Secundário). \\
\textbf{Nível de Meta} & Nível Usuário (\textit{Sea Level}). \\
\textbf{Pré-condições} & Auditoria em andamento pelo estudante e prazo da atividade vigente. \\
\textbf{Gatilho} & O estudante clica no botão "Finalizar e Submeter Auditoria" ou o cronômetro atinge o tempo zero (20 min). \\
\textbf{Fluxo Principal} & 
1. O sistema efetua a verificação de integridade estrutural das seções.\\
& 2. O sistema consolida o array \texttt{identified\_triggers} e o objeto \texttt{quality\_tools\_data}.\\
& 3. O sistema abre modal de confirmação irrevogável de envio.\\
& 4. O estudante confirma o envio da atividade.\\
& 5. O backend executa transação ACID atômica gravando na tabela \texttt{activity\_submissions}: \texttt{activity\_id}, \texttt{student\_id}, os dados JSONB e carimbo temporal \texttt{submission\_date}.\\
& 6. O sistema cancela o cronômetro regressivo da tela.\\
& 7. O sistema exibe notificação de envio bem-sucedido e bloqueia novas edições. \\
\textbf{Fluxos Alternativos} & 
\textbf{[4a] Submissão Forçada por Tempo}: Ao expirar os vinte minutos controlados por \texttt{SYSTEM\_TIMER} em \texttt{UC12}, os passos 3 e 4 são ignorados e o sistema dispara a persistência imediata no backend com os dados registrados até aquele segundo exato. \\
\textbf{Fluxos de Exceção} & 
\textbf{[5e] Falha Transacional de Banco}: Caso ocorra desconexão de rede ou colisão de integridade, a transação é revertida integralmente (\textit{rollback}) e a mensagem de erro é exibida para retentativa discente. \\
\textbf{Pós-condições} & Auditoria salva no banco de dados com nota nula aguardando avaliação docente. \\
\textbf{Regras de Negócio} & \textbf{RN12} (Submissão Única por Discente por Atividade) e Atomicidade ACID. \\
\textbf{Requisitos Associados} & \texttt{RF21}, \texttt{RNF01}, \texttt{RNF02}. \\
\bottomrule
\end{tabularx}
\fonte{Elaborado pelo autor (2026).}
\end{table}
